\documentclass[10pt,a4paper]{amsart}
\usepackage[margin=19mm]{geometry}
\usepackage{amsmath,amssymb,mathtools}
\usepackage[scr=rsfs,cal=euler]{mathalfa}
\usepackage{xfrac}
\usepackage[unicode,hidelinks,bookmarks=false]{hyperref}
\newcommand{\ie}{i.e.\ }
\newcommand{\cf}{cf.\ }
\newcommand{\resp}{respectively\ }
\newcommand{\Subsection}{\subsection}
\providecommand{\nobreakdash}{\nobreak\hbox{-}}
\setlength{\emergencystretch}{2em}
\multlinegap0pt
\usepackage{bm}
\usepackage{bbm}
\usepackage{dsfont}
\usepackage{xspace}
\usepackage{cancel}
\usepackage{mathtools}
\usepackage[shortlabels]{enumitem}
\usepackage{amsthm}
\makeatletter
\@ifundefined{theorem}{\newtheorem{theorem}{Theorem}}{}
\@ifundefined{lemma}{\newtheorem{lemma}[theorem]{Lemma}}{}
\makeatother
\DeclareMathOperator{\CC}{\mathbb C}
\DeclareMathOperator{\NN}{\mathbb N}
\DeclareMathOperator{\Span}{Span}
\DeclareMathOperator{\ZZ}{\mathbb Z}
\title[Computation of genus 2 Kleinian hyperelliptic functions via ]{Computation of genus 2 Kleinian hyperelliptic functions via Richelot isogenies}
\author{Matvey Smirnov}
\date{Source: arXiv:2603.23188v1 (CC BY 4.0)}
\begin{document}
\maketitle

\noindent\emph{Source and licence.} Computation of genus 2 Kleinian hyperelliptic functions via Richelot isogenies. Version arXiv:2603.23188v1 (CC BY 4.0), distributed under the Creative Commons Attribution 4.0 International licence, as recorded in the arXiv licence metadata for this exact version and shown on the abstract page https://arxiv.org/abs/2603.23188v1. Full rights notice: licenses/arxiv-2603.23188-CC-BY-4.0.txt.\par\smallskip

\noindent\textit{Editorial scope.} Counted unit: the Lemma whose source label is \texttt{lemThetaConvergence} in the article (source0.tex lines 632--638, source label \texttt{lemThetaConvergence}), delivered together with the article's own complete original proof of it (source0.tex lines 639--704, one \texttt{proof} environment of 66 lines). No mathematics is written, repaired or re-derived: statement and proof are the article's own text line for line. Required context retained uncounted: none. Every definition, display and label of those lines is retained unchanged. Omitted from this package and not redistributed: 1--631, 705--1001. Nothing inside a retained excerpt is omitted or altered: every retained body is delivered exactly as the article's lines are written, so no omission receipt of any kind accompanies this package. Citations: the retained excerpts carry no citation command at all, so no bibliography entry of the article is transcribed and no reference is redistributed. Reference adaptation: none was needed and none was made; every reference inside the delivered unit resolves inside it, so no unresolved reference or citation is produced. Printed slips kept literal: no printed word, symbol, hypothesis or number of the counted statement or of its proof was altered. Layout and class: the article's own class is \texttt{amsart} with packages including euscript, amsmath, amssymb, amsfonts, mathtools, mathrsfs, mathabx, stmaryrd, crossreftools, hyperref, tikz-cd, enumitem, graphicx, icomma; the wrapper is the lane's pinned \texttt{amsart} editorial wrapper at 10pt on A4, which restates the theorem-like environments the retained text needs and adds the article's own macro definitions that the retained text needs (\texttt{\textbackslash CC}, \texttt{\textbackslash NN}, \texttt{\textbackslash Span}, \texttt{\textbackslash ZZ}), with extra wrapper packages (bm, bbm, dsfont, xspace, cancel, mathtools, enumitem). The delivered numbering is the unit's own. Assets: no retained span contains a figure environment, a graphics-inclusion command, a PDF-page inclusion, a table or a TikZ picture; no image file is copied into this package, the package declares no asset dependency and no image file is redistributed with it. Licence: version arXiv:2603.23188v1 of this article is distributed under the Creative Commons Attribution 4.0 International licence, as recorded in the arXiv licence metadata for this exact version and shown on the abstract page https://arxiv.org/abs/2603.23188v1. A full rights notice is delivered at licenses/arxiv-2603.23188-CC-BY-4.0.txt. Characters: every retained excerpt is pure ASCII, so the delivered engine, XeLaTeX with its default Latin Modern text font, reports no missing character.
	\begin{lemma}\label{lemThetaConvergence}
		Assume that $\Omega$ is a quasi-reduced Riemann matrix and let $\{c_n\}_{n \in \mathbb N}$ be any sequence of positive real numbers that converges to $+\infty$. Consider any sequence $\{f_n\}_{n \in \NN}$, where $f_n \in R_2^{c_n \Omega}$. Then the following statements hold.
		\begin{enumerate}[label=(\roman*)]
			\item\label{thetaconvergencei} The sequence $\{f_n\}$ converges uniformly on compact subsets in $\CC^2$ if only if the sequence $p_n$ converges, where $p_n$ is the order $2$ Taylor expansion of $f_n$ at zero (i.e. $f_n(z) = p_n(z) + \bar{o}(z^2)$.
			\item\label{thetaconvergenceii} If the sequence $\{f_n\}$ converges uniformly on compact subsets in $\CC^2$, and $f$ denotes the limit of this sequence, then \[f \in \Span\{1, \sin^2(\pi z_1), \sin^2(\pi z_2),\sin^2(\pi(z_1 - z_2))\}.\]
		\end{enumerate}
	\end{lemma}

	\begin{proof}
		We shall denote Fourier coefficients of a function $g$ by $g\langle\cdot\rangle$. From the definition of the space $R_2^\Omega$ is is easy to see that a function $g \in R_2^\Omega$ can be expressed as a Fourier series
		\[
		g(z) = \sum_{m \in \ZZ^2}g\langle m\rangle \exp\left(2\pi i m^Tz\right),
		\]
		where the coefficients satisfy the relations
		\begin{equation}\label{eqFourierRecursion}
			g\langle m + 2k\rangle = \exp\left(2\pi i(m+ 2k)^T\Omega(m+ 2k) - 2\pi im^T\Omega m\right)
		\end{equation}
		for all $m,k \in \ZZ^2$. With this preparation we can prove an auxiliary fact: the sequence $\{f_n\}_{n \in \NN}$ converges uniformly on compact subsets of $\CC^2$ if and only if the sequences of Fourier coefficients $\left\{f_n\langle k\rangle \right\}_{n \in \NN}$ converge for all $k \in E \cup \{0\}$. Indeed, if the sequence $\{f_n\}_{n \in \NN}$ converges, then all Fourier coefficients are convergent sequences, since Fourier coefficient can be restored from the function by integration. On the other hand, assume that the sequences $\left\{f_n\langle k\rangle \right\}_{n \in \NN}$ converge for all $k \in E \cup \{0\}$. From~\eqref{eqFourierRecursion} it is evident that
		\begin{multline}\label{eqFourierExpansion}
			f_n(z) = f_n\langle 0\rangle\sum_{k \in \ZZ^2}\exp\left(8\pi ic_nk^T\Omega k\right)\exp\left(2\pi i (2k)^Tz\right) + \\ \sum_{m \in E} f_n\langle m\rangle \sum_{k \in \ZZ^2}\exp\left[2\pi ic_n\left((m+ 2k)^T\Omega(m+ 2k) - m^T\Omega m\right)\right]\exp(2\pi i(2k + m)^Tz).
		\end{multline}
		It is easy to see that the condition $c_n \to +\infty$ implies that when $n \to +\infty$ we have
		\[
		\sum_{k \in \ZZ^2}\exp\left(8\pi ic_nk^T\Omega k\right)\exp\left(2\pi i (2k)^Tz\right) \to 1
		\]
		uniformly with respect to $z$ on compact sets in $\CC^2$. Using the fact that $\Omega$ is quasi-reduced we also find that
		\begin{multline*}
			\sum_{k \in \ZZ^2}\exp\left[2\pi ic_n\left((m+ 2k)^T\Omega(m+ 2k) - m^T\Omega m\right)\right]\exp(2\pi i(2k + m)^Tz) \to \\ \exp(2\pi i m^Tz) + \exp(-2\pi i m^Tz) = 2\cos(2\pi m^tz)
		\end{multline*}
		uniformly on compact sets in $\CC^2$ (that is, the only terms that survive passing to the limit correspond to $k = 0$ and $k = -m$). Thus, the auxiliary statement is proved. In the meantime we also have found a formula for the limit, namely
		\[
		\lim_{n \to +\infty} f_n(z) = \lim_{n \to +\infty} f\langle 0\rangle  + \sum_{m \in E}\left(\lim_{n \to +\infty} f_n\langle m \rangle\right)2\cos(2\pi m^tz).
		\]
		The statement~\ref{thetaconvergenceii} is, therefore, proved.
		
		Now we prove~\ref{thetaconvergencei}. Clearly, if the sequence $\{f_n\}_{n \in \NN}$ converges uniformly in a neighbourhood of zero, then the Taylor expansions also constitute a convergent sequence. Thus, it remains to prove the converse. Let $p_n = p_n^{(0)} + p_n^{(11)}z_1^2 + p_n^{(12)}z_1z_2 + p_n^{(22)}z_2^2$ and assume that all the coefficients converge when $n \to +\infty$. From~\eqref{eqFourierExpansion} it is easy to calculate coefficients of $p_n$ by using Fourier coefficients $f_n \langle 0\rangle$ and $f_n\langle k\rangle$, $k \in E$. Indeed, if we put
		\[
		\alpha_1 = \begin{pmatrix}
			1 & 0
		\end{pmatrix}^T,\;\;\alpha_2 = \begin{pmatrix}
			0 & 1
		\end{pmatrix}^T,\;\;\alpha_3 = \begin{pmatrix}
			1 & -1
		\end{pmatrix}^T,
		\]
		then we can calculate that
		\[
		\begin{pmatrix}
			p_n^{(0)} \\ p_n^{(11)} \\ p_n^{(12)} \\ p_n^{(22)}
		\end{pmatrix} = A_n\begin{pmatrix}
			f_n\langle 0\rangle \\ f_n\langle \alpha_1 \rangle \\
			f_n\langle \alpha_2 \rangle \\ f_n\langle \alpha_3 \rangle
		\end{pmatrix}
		\]
		with a suitable matrix $A_n$ (which is universal for all functions in $R_2^{c_n\Omega}$). The formula for entries of $A_n$ is quite easy to obtain but we shall not give it here, since all that we are interested in is the limit of $A_n$ when $n \to +\infty$, which is easily computed as
		\[
		\lim_{n \to +\infty} A_n = \begin{pmatrix}
			1 & 2 & 2 & 2 \\
			0 & -4\pi^2 & 0 & -4\pi^2\\
			0 & 0 & 0 & 8\pi^2\\
			0 & 0 & -4\pi^2 & -4\pi^2
		\end{pmatrix}.
		\]
		Thus, we conclude that the limit of $A_n$ is an invertible matrix, so for large $n$ all matrices $A_n$ are invertible and, therefore, for large $n$ we can write that \[
		\begin{pmatrix}
			f_n\langle 0\rangle \\ f_n\langle \alpha_1 \rangle \\
			f_n\langle \alpha_2 \rangle \\ f_n\langle \alpha_3 \rangle
		\end{pmatrix} = A_n^{-1}
		\begin{pmatrix}
			p_n^{(0)} \\ p_n^{(11)} \\ p_n^{(12)} \\ p_n^{(22)}
		\end{pmatrix}
		\]
		Since on the right-hand side of this equation all sequences converge (note that the matrix inversion is continuous), we obtain that Fourier coefficients of functions $f_n$ converge when $n \to +\infty$. Thus, using our auxiliary statement we get~\ref{thetaconvergencei}.
	\end{proof}

\end{document}
